\documentclass[11pt]{article}
\usepackage[margin=0.85in]{geometry}
\usepackage[T1]{fontenc}
\usepackage[utf8]{inputenc}
\usepackage{lmodern}
\usepackage{microtype}
\usepackage{amsmath}
\usepackage{booktabs}
\usepackage{tabularx}
\usepackage{enumitem}
\usepackage{titlesec}
\usepackage[hidelinks]{hyperref}

\setlist{nosep, leftmargin=*}
\titlespacing*{\section}{0pt}{1.0ex}{0.5ex}
\setlength{\parskip}{0.5ex}
\setlength{\parindent}{0pt}
\newcolumntype{Y}{>{\raggedright\arraybackslash}X}
\newcommand{\rhob}{\bar{\rho}}

\begin{document}

\begin{center}
  {\Large\bfseries Lab Report 1 (Weeks 1--2): Agents for a Lattice-Stiffness GNN}\\[0.5ex]
  Jamie Carpenter \quad$\cdot$\quad AM/ME/CE 295 \quad$\cdot$\quad 18 October 2026 \quad$\cdot$\quad
  \url{https://github.com/jamie-carpenter/lattice-gnn}
\end{center}

\section{Draft project proposal}

\textbf{Problem and goal.} The effective elastic properties of architected strut lattices usually
come from finite-element (FE) homogenization, which is accurate but too slow for exploring large
design spaces of unit cells. This project builds a graph neural network (GNN) that predicts these
properties directly from the unit-cell geometry and generalizes to unit cells it has not seen in
training. Each unit cell is represented as a periodic graph, with nodes at strut junctions and edges
along struts, and the model predicts $\log_{10}$ of the moduli $E_x, E_y, E_z, G_{yz}, G_{xz},
G_{xy}$ (normalized by the base-material modulus $E_s$) and three Poisson coupling terms. It is
trained on the catalog of Lumpe \& Stankovi\'c~\cite{lumpe}, which contains about 17,000
lattices. The same framework could later be extended to manufacturing defects, such as
strut-diameter variation, waviness, missing struts and node misplacement, and possibly to size
effects, but these are future directions rather than the main focus. The intended users are
researchers and engineers who need fast, physically valid stiffness estimates for many candidate
unit cells.

\textbf{Why an agent-based approach.} Building a reliable surrogate is a chain of distinct,
tool-heavy sub-goals, each of which can be handled by its own agent while I approve every
consequential action. The non-agent baseline is the same pipeline done by hand, with a fixed
training script and a manual hyperparameter search.

\textbf{Planned architecture.} Table~\ref{tab:agents} breaks the goal into sub-goals and assigns
each one to its own agent, listing the data that agent is given (its memory) and the environment
it works in (its tools). All agents run in Claude Code with pinned model IDs and
share memory through versioned files in the repository, such as \texttt{LOG.md},
\texttt{papers.json} and the training-run records.

\begin{table}[h]\centering\small
\caption{Sub-goals, agents, memory and tools.}\label{tab:agents}
\begin{tabularx}{\linewidth}{@{}>{\raggedright\arraybackslash}p{2.9cm} >{\raggedright\arraybackslash}p{2.0cm} Y Y@{}}
\toprule
Sub-goal & Agent & Data supplied (memory) & Environment (tools) \\
\midrule
Understand conventions and prior work & Literature (built) & Project brief; writes \texttt{papers.json} and a review with page-cited evidence that the other agents read & arXiv, OpenAlex, citation links and a page-by-page PDF reader (a tool server), web search \\
Prepare and validate the data & Data & Catalog file, data-check report, fixed train/test splits (random, and by held-out lattice family) & Parser, periodic graph builder, \texttt{check\_data.py}, \texttt{plot\_lattice.py} \\
Design and train the GNN & Training & Conventions from the review, data report, experiment log of configurations and metrics & Python with PyTorch Geometric on the AWS WorkSpace, \texttt{train.py} \\
Tune and stress-test generalization & Evaluation & Experiment log, held-out unseen unit cells, recorded failure cases & AWS WorkSpace, \texttt{predict.py}, positive-definiteness check \\
Model defects (later) & Defects & Defect literature, defect specifications, database of FE results & FE solver (Abaqus or CalculiX) exposed as a tool server \\
\bottomrule
\end{tabularx}
\end{table}

\textbf{Measuring success and risks.} The main measure is the prediction error (RMSE and $R^2$) on
lattices held out from training. Many lattices in the catalog are close variants of one another,
so if the test set is drawn at random, most test lattices have a near-twin in the training set and
the error looks better than it would on a truly new design. I will therefore also hold out whole
families of related lattices, so that the test lattices have topologies the model has never seen,
and compare both results with the baselines. I will also assemble the predicted $6\times6$
stiffness matrix and count how often it is not positive definite, and every reported number must be
reproducible from a committed command. The main risks are unconfirmed dataset conventions
(Section~\ref{sec:rw}) and the FE pipeline needed for the defect extension.

\section{Related-work note}\label{sec:rw}

The literature agent recorded 22 papers. The three most relevant to this project are summarized
below. Page numbers refer to PDF pages of the arXiv versions, and I checked every number cited here
against the source.

\textbf{Lumpe \& Stankovi\'c, PNAS 2021~\cite{lumpe}: the dataset.} The authors decode crystal nets
into periodic strut lattices and compute their effective elastic properties with an
Euler--Bernoulli FE beam model. The base material has $E_s = 1$\,MPa and $\nu_s = 0.3$, and every
structure uses a constant bar radius of 0.1\,mm, scaled to a relative density of $\rhob = 0.01$.
After removing duplicates and structures with very small members, the catalog has 17,087 unique
entries (Methods, read as Europe PMC text without page numbers). This paper defines my training
labels and confirms my own finding that they correspond to 1\% relative density: solving
$E/E_s = C\rhob^{\,n}$ for each record gives a median of 1.00\%. My copy of the catalog has 17,262
records, 41 of them without properties, which I still need to reconcile with the paper's count.

\textbf{Indurkar et al.\ 2022~\cite{indurkar}: a GNN on the same catalog.} This work trains a
message-passing GNN on 17,201 lattices from the same catalog at $\rhob = 0.01$ [p.~1]. The graph
includes periodic edges to neighbouring unit cells [p.~3], the same design choice I use. Instead of
predicting stiffness, the network classifies lattices as stretching-dominated, bending-dominated or
mixed, with 81.8\% validation accuracy on a random 90:10 split [p.~4]. The moduli span $10^{-8}$ to
$10^{-3}$ [p.~2], which supports predicting $\log_{10}$ moduli. Their random split can put
near-identical lattices in both the training and the validation set, which is why I will also test
on held-out families of lattices.

\textbf{Grega et al., ICLR 2024~\cite{grega}: an equivariant GNN for the stiffness tensor.} This
work uses the same catalog [p.~4], cleaned to 8,954 base lattices and augmented with nodal
perturbations to 635,454 geometries, with the perturbations of each lattice kept on one side of the
train/test split [p.~5]. An equivariant MACE network predicts the full stiffness tensor; without
constraints it predicted physically impossible negative eigenvalues for 30\% of lattices, which a
positive-semi-definite output layer removed [p.~7]. This is the closest prior work to mine. It shows
that physical validity must be checked explicitly, and its nodal perturbations are a useful
template for the later defect extension.

\textbf{Synthesis.} Together, these papers show that the catalog is a sound basis for learning,
that periodic graphs are an effective representation, and that GNNs can predict lattice stiffness.
They also point to two pitfalls: test scores that look too good because near-identical lattices
appear in both the training and test sets, and predictions that are physically invalid. My project
builds on these papers, and I will evaluate my model on lattice families it has never seen, so that
its accuracy reflects how it would perform on genuinely new unit cells.

\section{Progress}\label{sec:progress}

\textbf{Done.} The repository contains the data pipeline (a parser, a periodic graph builder and
the GNN model in \texttt{lattice\_gnn/}), the full catalog, pinned requirements, and a running log
in \texttt{LOG.md}; the data check passes on the sample file. The literature agent, first built on
the Claude Agent SDK and then moved to Claude Code, was run with
\texttt{python -m agents.literature --name main}. The first run stopped after 32 turns when Claude
Code's shell-safety check refused a tool call containing a formula, and the agent stopped rather
than work around it; after I adjusted its instructions, the second run completed in 81 turns with a
page-cited review of 22 papers. On the course's Amazon WorkSpaces desktop, GitHub access, Python
3.12 with the pinned packages, Claude Code and network access all work, so code-running agents can
run there, isolated from my own computer and with my approval of each command. Abaqus is not yet
installed (full results in \texttt{LOG.md}).

\textbf{Remaining.} I still need to read the dataset's supplementary information, reconcile the
record counts, run the baselines and a first training run, and read the defect literature in full.

\textbf{Bottleneck.} The one open bottleneck is Abaqus: it is not installed on the WorkSpace, and I
have not yet confirmed that the Caltech licence allows its use on AWS. It only affects the later
defect work, so it does not block Lab Report 2. For now I can prototype defect labels with an
open-source solver (CalculiX or OpenSeesPy), and I will confirm the licence terms with Caltech IT
by 1 November.

\section{Next steps}

\begin{enumerate}
  \item Define the research question precisely and write the one-page project proposal.
  \item Read the dataset's supplementary information, reconcile the record counts, and update the
        related-work note.
  \item Implement the baselines and complete a first GNN training run on the CPU, testing on
        held-out lattice families and checking positive definiteness, and record the results in \texttt{LOG.md}.
  \item Confirm the Abaqus licence terms for AWS, choose the FE solver for the later defect work,
        and test it on a single lattice.
\end{enumerate}

\section{AI use note}

Claude Code (\texttt{claude-opus-5-5}) ran the environment checks, built the literature agent's
command-line tool interface, checked citations and drafted this report. The literature agent
(\texttt{claude-sonnet-5-5}, Claude Code 2.1.299) searched for, read and recorded the papers; its
outputs and run record are in \texttt{literature\_review/main/}. I verified the environment
results, chose the three papers, and checked every number and page citation in
Section~\ref{sec:rw} against the PDFs and the Europe PMC text. One citation from the agent turned
out to be a paraphrase rather than a direct quote, so I did not use it.

\begin{thebibliography}{3}\small
\bibitem{lumpe} T.~S. Lumpe, T.~Stankovi\'c. Exploring the property space of periodic cellular
  structures based on crystal networks. \emph{PNAS} 118(7), 2021. doi:10.1073/pnas.2003504118.
\bibitem{indurkar} P.~P. Indurkar, S.~Karlapati, A.~J.~D. Shaikeea, V.~S. Deshpande. Predicting
  deformation mechanisms in architected metamaterials using GNN. arXiv:2202.09427v2, 2022.
\bibitem{grega} I.~Grega, I.~Batatia, G.~Cs\'anyi, S.~Karlapati, et al. Energy-conserving
  equivariant GNN for elasticity of lattice architected metamaterials. ICLR 2024,
  arXiv:2401.16914v2.
\end{thebibliography}

\end{document}
